% BioNetFlux Workflow Section for Main Documentation
% To be included in master LaTeX document

\section{BioNetFlux Experimental Workflow}
\label{sec:bionetflux_workflow}

This section illustrates the typical experimental workflow in BioNetFlux, showing how the configuration-driven architecture coordinates the various framework components to perform biological network transport simulations.

\subsection{Workflow Overview}

The BioNetFlux framework follows a systematic six-phase approach for conducting biological transport simulations:

\begin{enumerate}
    \item \textbf{Configuration Phase}: Load and validate TOML configuration files
    \item \textbf{Geometry Phase}: Create and validate network topology
    \item \textbf{Problem Setup Phase}: Generate problem instances and discretizations
    \item \textbf{Solver Setup Phase}: Initialize numerical solution components
    \item \textbf{Time Evolution Phase}: Perform time-stepping with Newton iteration
    \item \textbf{Output Phase}: Generate visualizations and analysis
\end{enumerate}

\subsection{Detailed Workflow Diagram}

The following diagram shows the complete experimental workflow with class relationships and method calls:

\begin{figure}[H]
\centering
\begin{tikzpicture}[scale=0.8, every node/.style={scale=0.8}]

% Define colors
\definecolor{configcolor}{RGB}{255, 165, 0}
\definecolor{geometrycolor}{RGB}{0, 128, 128}
\definecolor{problemcolor}{RGB}{70, 130, 180}
\definecolor{solvercolor}{RGB}{220, 20, 60}
\definecolor{outputcolor}{RGB}{34, 139, 34}

% Define styles
\tikzset{
    configbox/.style={rectangle, rounded corners, minimum width=2cm, minimum height=0.8cm, 
                     align=center, draw=configcolor, fill=configcolor!20, thick, font=\scriptsize},
    geometrybox/.style={rectangle, rounded corners, minimum width=2cm, minimum height=0.8cm,
                       align=center, draw=geometrycolor, fill=geometrycolor!20, thick, font=\scriptsize},
    problembox/.style={rectangle, rounded corners, minimum width=2cm, minimum height=0.8cm,
                     align=center, draw=problemcolor, fill=problemcolor!20, thick, font=\scriptsize},
    solverbox/.style={rectangle, rounded corners, minimum width=2cm, minimum height=0.8cm,
                    align=center, draw=solvercolor, fill=solvercolor!20, thick, font=\scriptsize},
    outputbox/.style={rectangle, rounded corners, minimum width=2cm, minimum height=0.8cm,
                     align=center, draw=outputcolor, fill=outputcolor!20, thick, font=\scriptsize},
    arrow/.style={->, >=stealth, thick}
}

% Configuration Phase
\node[configbox] (config) {TOML\\ Config};
\node[configbox, right=1.5cm of config] (configmgr) {Config\\Manager};
\node[configbox, right=1.5cm of configmgr] (resolver) {Function\\Resolver};

% Geometry Phase  
\node[geometrybox, below=1cm of config] (geometry) {Domain\\Geometry};
\node[geometrybox, right=1.5cm of geometry] (factory) {Geometry\\Factory};

% Problem Phase
\node[problembox, below=1cm of geometry] (problems) {Problem\\List};
\node[problembox, right=1.5cm of problems] (discretization) {Discretization\\List};
\node[problembox, right=1.5cm of discretization] (constraints) {Constraint\\Manager};

% Solver Phase
\node[solverbox, below=1cm of problems] (global_disc) {Global\\Discretization};
\node[solverbox, right=1.5cm of global_disc] (assembler) {Global\\Assembler};

% Time Evolution
\node[solverbox, below=1cm of global_disc] (newton) {Newton\\Solver};
\node[solverbox, right=1.5cm of newton] (timestepper) {Time\\Stepper};

% Output
\node[outputbox, below=1cm of newton] (plotter) {Plotter};
\node[outputbox, right=1.5cm of plotter] (analysis) {Analysis};

% Arrows showing data flow
\draw[arrow, configcolor] (config) -- (configmgr);
\draw[arrow, configcolor] (configmgr) -- (resolver);
\draw[arrow, geometrycolor] (configmgr) -- (geometry);
\draw[arrow, geometrycolor] (geometry) -- (factory);
\draw[arrow, problemcolor] (geometry) -- (problems);
\draw[arrow, problemcolor] (geometry) -- (discretization);
\draw[arrow, problemcolor] (geometry) -- (constraints);
\draw[arrow, problemcolor] (resolver) -- (problems);
\draw[arrow, solvercolor] (discretization) -- (global_disc);
\draw[arrow, solvercolor] (problems) -- (assembler);
\draw[arrow, solvercolor] (constraints) -- (assembler);
\draw[arrow, solvercolor] (global_disc) -- (assembler);
\draw[arrow, solvercolor] (assembler) -- (newton);
\draw[arrow, solvercolor] (newton) -- (timestepper);
\draw[arrow, outputcolor] (timestepper) -- (plotter);
\draw[arrow, outputcolor] (timestepper) -- (analysis);

% Phase labels
\node[above=0.5cm of configmgr, configcolor, font=\bfseries\small] {Configuration};
\node[above=0.5cm of factory, geometrycolor, font=\bfseries\small] {Geometry};
\node[above=0.5cm of discretization, problemcolor, font=\bfseries\small] {Problem Setup};
\node[above=0.5cm of assembler, solvercolor, font=\bfseries\small] {Solver Setup};
\node[below=0.3cm of timestepper, orange, font=\bfseries\small] {Time Evolution};
\node[below=0.3cm of analysis, outputcolor, font=\bfseries\small] {Output};

\end{tikzpicture}
\caption{BioNetFlux Experimental Workflow - Class Relationships and Data Flow}
\label{fig:bionetflux_workflow}
\end{figure}

\subsection{Key Workflow Steps}

\subsubsection{Phase 1: Configuration Loading}

\begin{lstlisting}[language=Python, caption=Configuration Phase Example]
# Load problem-specific configuration
from bionetflux.problems.ooc_config_manager import OoCConfigManager

config_manager = OoCConfigManager()
config = config_manager.load_config("my_experiment.toml")

# Functions are automatically resolved from string specifications
initial_conditions = config['initial_conditions']  # Now callables
force_functions = config['force_functions']        # Now callables
\end{lstlisting}

\subsubsection{Phase 2: Geometry Creation}

\begin{lstlisting}[language=Python, caption=Geometry Phase Example]
# Create or load geometry
from bionetflux.geometry.domain_geometry import build_grid_geometry

geometry = build_grid_geometry(N=4)  # 4×4 grid with 18 domains
geometry.validate_geometry(verbose=True)

# Geometry automatically tracks connections
boundary_connections = geometry.get_boundary_connections()
interior_connections = geometry.get_interior_connections()
\end{lstlisting}

\subsubsection{Phase 3: Problem Setup}

\begin{lstlisting}[language=Python, caption=Problem Setup Phase Example]
# Create problems from geometry and configuration
problems = []
for domain_id in range(geometry.num_domains()):
    domain_info = geometry.get_domain(domain_id)
    
    problem = Problem(
        neq=config['problem']['neq'],
        domain_start=domain_info.domain_start,
        domain_length=domain_info.domain_length,
        parameters=parameters_from_config,
        problem_type="organ_on_chip"
    )
    
    # Apply configuration functions
    problem.set_initial_condition(0, initial_conditions['u'])
    problem.set_force(0, force_functions['u'])
    problems.append(problem)
\end{lstlisting}

\subsubsection{Phase 4: Solver Setup}

\begin{lstlisting}[language=Python, caption=Solver Setup Phase Example]
# Create global discretization and assembler
global_disc = GlobalDiscretization(discretizations)
global_disc.set_time_parameters(dt, T)

# Setup constraints from geometry
constraint_manager = setup_constraints_from_geometry(geometry, problems, neq)
constraint_manager.map_to_discretizations(discretizations)

# Initialize global assembler
assembler = GlobalAssembler(problems, global_disc, constraint_manager)
\end{lstlisting}

\subsubsection{Phase 5: Time Evolution}

\begin{lstlisting}[language=Python, caption=Time Evolution Phase Example]
# Main simulation loop
newton_solver = NewtonSolver(tolerance=1e-10)
time_stepper = TimeStepper(dt, T)

while not time_stepper.is_finished():
    # Assemble system for current time
    residual, jacobian = assembler.assemble_residual_and_jacobian(
        current_solution, time_stepper.current_time
    )
    
    # Solve nonlinear system
    solution_update = newton_solver.solve(residual, jacobian)
    
    # Advance time
    current_solution = time_stepper.advance_time(
        current_solution, solution_update
    )
\end{lstlisting}

\subsubsection{Phase 6: Output and Analysis}

\begin{lstlisting}[language=Python, caption=Output Phase Example]
# Extract domain solutions
domain_solutions = assembler.extract_domain_solutions(final_solution)

# Visualization
plotter = LeanMatplotlibPlotter(problems, discretizations)
plotter.plot_2d_curves(domain_solutions, title="Final Solution")
plotter.plot_birdview(domain_solutions, equation_idx=0, time=T)
plotter.plot_flat_3d(domain_solutions, equation_idx=0)

# Analysis
total_mass = compute_total_mass(domain_solutions)
conservation_error = check_mass_conservation(domain_solutions)
\end{lstlisting}

\subsection{Complete Experiment Example}

Here's a complete example showing how all phases work together:

\begin{lstlisting}[language=Python, caption=Complete BioNetFlux Experiment]
def run_bionetflux_experiment(config_file="ooc_experiment.toml"):
    """
    Complete BioNetFlux experimental workflow.
    """
    
    # Phase 1: Configuration
    from bionetflux.problems.ooc_problem import create_global_framework
    problems, global_disc, constraints, name = create_global_framework(
        config_file=config_file
    )
    
    # Phase 2-4: Setup (handled by create_global_framework)
    print(f"Experiment: {name}")
    print(f"Domains: {len(problems)}")
    print(f"Constraints: {constraints.n_constraints}")
    
    # Phase 5: Time Evolution
    from bionetflux.solvers.setup_solver import setup_complete_solver
    solver = setup_complete_solver(problems, global_disc, constraints)
    
    # Create initial conditions
    initial_traces, initial_multipliers = solver.create_initial_conditions()
    
    # Time evolution loop
    final_solution = solver.evolve_in_time(
        initial_solution=(initial_traces, initial_multipliers),
        progress_callback=lambda t, sol: print(f"Time: {t:.3f}")
    )
    
    # Phase 6: Output
    from bionetflux.visualization.lean_matplotlib_plotter import LeanMatplotlibPlotter
    
    plotter = LeanMatplotlibPlotter(problems, global_disc.spatial_discretizations)
    final_traces, _ = solver.extract_domain_solutions(final_solution)
    
    # Generate comprehensive output
    plotter.plot_2d_curves(final_traces, title=f"{name} - Final State")
    plotter.plot_comparison(initial_traces, final_traces, 0.0, global_disc.T)
    plotter.show_all()
    
    return final_traces, solver

# Run experiment
if __name__ == "__main__":
    results, solver = run_bionetflux_experiment("my_ooc_experiment.toml")
    print("Experiment completed successfully!")
\end{lstlisting}

\subsection{Workflow Benefits}

The configuration-driven workflow provides several advantages:

\begin{itemize}
    \item \textbf{Reproducibility}: TOML files ensure exact experimental reproduction
    \item \textbf{Flexibility}: Easy parameter sweeps and function variations
    \item \textbf{Modularity}: Clear separation between setup, solving, and analysis
    \item \textbf{Validation}: Automatic consistency checking at each phase
    \item \textbf{Extensibility}: New problem types integrate seamlessly
    \item \textbf{Documentation}: Configuration files serve as experimental records
\end{itemize}

This workflow architecture positions BioNetFlux as a powerful platform for systematic biological network research with excellent experimental control and documentation capabilities.
